%%%%%%%%%
\section{Graded unitarity constraints on $V_k(\mathfrak{so}_{2N})$} \label{sec:typeDgrunit}
%%%%%%%%%
In this section, we investigate constraints imposed by graded unitarity on $V_k(\mathfrak{so}_{2N})$ affine Kac--Moody algebras with $N \geqslant 4$. The strategy is similar to the previous section, however, as we elaborate in this section, the analysis is technically complicated due to the intricate structure of zeros of the Gorelik--Kac vacuum determinant. 

We begin by recalling that the positive roots and the Weyl vector for $\sof_{2N}$ are given by,
\begin{equation}
\begin{split}
    &\Delta_+ = \{\alpha^\pm_{r,s} := e_r \pm e_s: 1 \leqslant r < s \leqslant N \}, \\
    &\rho = \frac{1}{2} \sum_{\alpha \in \Delta_+} \alpha = \left(N-1, N-2, \ldots, 1,0 \right),
\end{split}
\end{equation}
with $(\rho|\alpha_{r,s}^+) = 2N - r - s$ and $(\rho|\alpha_{r,s}^-) = s - r$.

As in the type $A$ affine Kac--Moody case, we first simplify the regularity criterion \eqref{eq:ADEregularity} for $\lambda = t \alpha$ where $t \in \mathbb{Z}_{>0}$ and $\alpha \in \Delta_+$. It can be shown by direct computation that for $\gf = \sof_{2N}$ and $\lambda = t \alpha^{\pm}_{r,s}$, the regularity criterion simplifies to,
\begin{equation} \label{eq:typeDregularity}
\begin{split}
     &E(t \alpha^+_{r,s}) \neq 0 \Longleftrightarrow \begin{cases}
         t \geqslant s \qquad \text{or} \\ 
         t \geqslant r \ \ \text{and} \ \ t = s - r
     \end{cases} \\
     &E(t \alpha^-_{r,s}) \neq 0 \Longleftrightarrow \begin{cases}
         t \geqslant r \ \ \text{and} \ \ t \geqslant 2N - s \qquad \text{or} \\ 
         t \geqslant r \ \ \text{and} \ \ t = 2N - 2s \qquad \text{or} \\ 
         t \geqslant r \ \ \text{and} \ \ t = 2N-r-s
     \end{cases}     
\end{split}
\end{equation}
Following a similar approach as in the previous section, we first determine the smallest values of $A_{(t,\alpha)} = t(t+(\rho|\alpha))$ consistent with the simplified regularity constraints in \eqref{eq:typeDregularity} to be,
\begin{equation} \label{eq:AtconsttypeD}
\begin{split}
    &A_{(t,\alpha^+)} \ \ \begin{cases}
        = (2N-2) \qquad \text{when} \ \ t = 1 \ \ \text{and} \ \ \alpha^+ = \alpha^+_{1,2} = \theta \\ \geqslant (4N-8) \qquad \text{otherwise}
    \end{cases} \\
    &A_{(t,\alpha^-)} \ \ \begin{cases}
        = (2N-2) \qquad \text{when} \ \ t = 2 \ \ \text{and} \ \ \alpha^- = \alpha^-_{2,N-1} \\
        = 2N \qquad \qquad \ \text{when} \ \ t = 2 \ \ \text{and} \ \ \alpha^- = \alpha^-_{1,N-1} \\ \geqslant (4N-8) \qquad \text{otherwise}
    \end{cases}
\end{split}
\end{equation}

As a direct consequence of \eqref{eq:AtconsttypeD}, it can be observed that $E(\theta)$ and $E(2\alpha^-_{2,N-1})$ are relevant for computing the multiplicity of the zero $k = k_{2N-2,L}$, while $E(2\alpha^-_{1,N-1})$ is relevant for computing multiplicity of the $k = k_{2N,L}$ zero in the conformal level $L$ vacuum determinant in a fixed-weight sector. The multiplicities of the zeros $k = k_{2N-2,L}$ and $k = k_{2N,L}$ can then be determined by tracking the possible occurrence of corresponding pairs $(t,\alpha)$ in the summation in \eqref{eq:splitM}. Thus using \eqref{eq:expinchar}, for $N \geqslant 4$, we find that the multiplicity of the $k = k_{2N-2,L}$ zero in the fixed root, conformal level $L$ vacuum determinant is given by,\footnote{We note that while the weights $2\alpha^-_{2,N-1}$ and $2\alpha^-_{1,N-1}$ are not dominant, it can be shown that $(s_N s_{N-1}s_1).(2\alpha^-_{2,N-1}) = \theta$ and $(s_N s_{N-1}).(2\alpha^-_{1,N-1}) = 2\omega_1$. Thus, $E(2\alpha^-_{2,N-1}) = -E(\theta)$ and $E(2\alpha^-_{1,N-1}) = E(2\omega_1)$.} 
\begin{equation} \label{eq:Lmultofk2Nm2}
\begin{split}
    &N \ \text{even}: \qquad m_{(2N-2)/L}(L\delta+\alpha) = \begin{cases}
        1 \qquad L \in 2\mathbb{N}+1 \ \ \text{and} \ \ L \neq 0~\text{mod}~(N-1) \\
        2 \qquad L \in 2\mathbb{N}+1 \ \ \text{and} \ \ L = 0~\text{mod}~(N-1) \\
        0 \qquad L \in 2\mathbb{N}
    \end{cases} \\
    &N \ \text{odd}: \qquad m_{(2N-2)/L}(L\delta+\alpha) = \begin{cases}
        1 \qquad L \in 2\mathbb{N}+1  \\
        1 \qquad L = (N-1)~\text{mod}~(2N-2) \\
        0 \qquad \text{otherwise}
    \end{cases}
\end{split}
\end{equation}
In the same manner, for $N \geqslant 5$, the multiplicity of the $k = k_{2N,L}$ zero in the fixed-root vacuum determinant at conformal level $L$ is given by,
\begin{equation} \label{eq:Lmultofk2N}
   m_{2N/L}(L\delta+\alpha) = \begin{cases}
        1 \qquad L \in 2\mathbb{N} \ \ \text{and} \ \ L \neq 0~\text{mod}~N \\
        0 \qquad \text{otherwise}
    \end{cases}
\end{equation}
To understand the structure of zeros of the all-level vacuum determinant, we note from \eqref{eq:AtconsttypeD}, \eqref{eq:Lmultofk2Nm2}, and \eqref{eq:Lmultofk2N}, that the (possible) zeros of the conformal level $L$ vacuum determinant in a fixed-root sector for $V_k(\sof_{2N})$ can be classified as,
\begin{equation} \label{eq: so2nprelimzero}
    k = \begin{cases}
        k_{2N-2,q} = -(2N-2) + \frac{2N-2}{q}: q \leqslant L, \ q \neq 0~\mathrm{mod}~(N-1), \ \ \text{and} \ \ q \ \mathrm{odd} \\ 
        k_{2N-2,q} = -(2N-2) + \frac{2N-2}{q}: q \leqslant L, \ q = (N-1)~\mathrm{mod}~(2N-2) \\
        k_{2N,q} = -(2N-2) + \frac{2N}{q} \qquad: q \leqslant L, \ q \neq 0~\mathrm{mod}~N, \ \ \text{and} \ \ q \ \mathrm{even} \\
        k_{A,B} = -(2N-2) + \frac{A}{B} \qquad \ \ : A \geqslant (4N-8) \ \ \text{and} \ \ B \leqslant L
    \end{cases}
\end{equation}
As a consequence of the multiplicities in \eqref{eq:Lmultofk2Nm2}, the structure of zeros and the constraints from graded unitarity are qualitatively different for even and odd $N$. To illustrate the general strategy, we begin with a detailed analysis of $V_k(\mathfrak{so}_8)$, before treating the general case of even $N$. The corresponding analysis for odd $N$ is deferred to Appendix \ref{app:VkSOodd}, and we conclude the section by summarising graded unitarity constraints for all $N \geqslant 5$. 

\subsection{$V_k(\sof_8)$}
Implementing graded unitarity constraints on $V_k(\sof_8)$ requires a meticulous study of multiplicities of zeros in the vacuum determinant in a fixed-weight sector, and the representations corresponding to the specific states becoming null at these zeros. We will first outline the structure of zeros of the fixed-root vacuum determinant at conformal levels $L \leqslant 6$, and then infer the constraints obtained by graded unitarity. This bespoke analysis will orient us for investigating all-level determinant-based graded unitarity constraints by invoking the sign constraint in \eqref{eq:grunitsignrule}.

\subsubsection{Low level graded unitarity constraints}
The vacuum determinant in a fixed-root sector for conformal levels $L = 1$ and $L = 2$ takes the following form, 
\begin{equation} \label{eq:lowlevelso8det1to2}
\begin{split}
    &\text{det}~S_{\delta + \alpha} \propto (k-k_{6,1}) \\
    &\qquad \qquad \propto \textcolor{blue}{(k - 0)}~, \\
    &\text{det}~S_{2\delta + \alpha} \propto (k - k_{4,1})^3(k-k_{6,1})^{\#_1} (\ldots \ldots)\\ 
    &\qquad \qquad  \propto \textcolor{blue}{(k + 2)^3}(k - 0)^{\#_1} (\ldots \ldots)~. \\ 
\end{split}
\end{equation}
Level one constraints combined with the requirement that $R(T) = 1$ immediately implies that $-6 < k \leqslant 0$. 
At $L = 2$, the multiplicity three of $(k+2)$ in a fixed-root sector is a manifestation of the triality symmetry of $\sof_8$. In particular, it can be shown that the zero at $k = -2$ is a consequence of three null states, which are components of the $[2,0,0,0]$, $[0,0,2,0]$ and $[0,0,0,2]$ representations. Consequently, one sees that the graded unitarity sign rule eliminates $k \in (-2,0)$, leading to the constraint $-6<k \leqslant -2$ or $k = 0$. 

The constraints from graded unitarity are more interesting and involved at $L \geqslant 3$. To this end, note that the vacuum determinant in a fixed-root sector at $L = 3,4,5,6$ is given by,
\begin{equation} \label{eq:lowlevelso8det3to6}
\begin{split}
    &\text{det}~S_{3\delta + \alpha} \propto (k - k_{2,1})^2 (k-k_{4,1})^{\#_2} (\ldots \ldots)\\ 
    &\qquad \qquad  \propto \textcolor{blue}{(k + 4)^2} (k +2 )^{\#_2} (\ldots \ldots)~,\\
    &\text{det}~S_{4\delta + \alpha} \propto (k - k_{2,1})^{32}(k-k_{7,2})^5(k-k_{4,1})^{\#_3} (\ldots \ldots)\\ 
    &\qquad \qquad  \propto \textcolor{blue}{(k + 4)^{32}(k + 5/2)^5}(k+2)^{\#_3} (\ldots \ldots)~, \\
    &\text{det}~S_{5\delta + \alpha} \propto (k - k_{6,5})(k-k_{2,1})^{322} (k-k_{7,2})^{\#_4} (\ldots \ldots)\\ 
    &\qquad \qquad  \propto \textcolor{blue}{(k+24/5)(k + 4)^{322}}(k + 5/2)^{\#_4} (\ldots \ldots)~, \\
    &\text{det}~S_{6\delta + \alpha} \propto (k - k_{6,5})^{\#_5}(k-k_{4,3})^3 (\ldots \ldots)\\ 
    &\qquad \qquad  \propto \textcolor{blue}{(k+24/5)^{\#_5}(k + 14/3)^3} (\ldots \ldots)~. \\
\end{split}
\end{equation}
First, note that at $L = 3$, there are two primitive null states at $k = -4$ and both are components of the adjoint representation, hence contributing to the multiplicity two in a fixed-root sector. As a result, \textit{even though} there is a \textit{new smallest zero} at $k = -4$ in the vacuum determinant at $L = 3$, there are no new fixed root determinant-based constraints from the graded unitarity sign rule \eqref{eq:grunitsignrule}. This is a first of many instances for $V_k(\sof_{2N})$ where the structure of zeros is qualitatively different from that encountered in the $V_k(\mathfrak{sl}_N)$ graded unitarity analysis. 

At $L = 4$, there is a single primitive null state at $k = -5/2$, which is a component of the $[0,2,0,0]$ representation. The multiplicity of five then follows from the root multiplicity of the $V_{2\omega_2}$ representation being five. Hence, at $L = 4$, graded unitarity constraints rule out $k \in (-5/2,-2)$. At $L = 5$, the primitive null state giving rise to the $k = -24/5$ zero is a component of the adjoint representation, and graded unitarity further rules out $k \in (-24/5,-4) \bigcup (-4, -5/2)$. At $L = 6$ there are no new constraints, and hence determinant-based graded unitarity constraints up to $L = 6$ restrict the level $k$ to be,
\begin{equation}
    k \leqslant -\frac{24}{5} \qquad \text{or} \qquad k = \left\{-4, -\frac{5}{2},-2,0 \right\}~.
\end{equation}
Extending these multiplicity computations to higher conformal levels, we find that fixed root determinant-based graded unitarity constraints up to $L = 12$ are given by,
\begin{equation} \label{eq:lowlevelso8fullconst}
    k \leqslant -\frac{60}{11} \qquad \text{or} \qquad k = \left\{-\frac{16}{3}, -\frac{26}{5}, -\frac{36}{7}, -\frac{24}{5}, -4, -\frac{5}{2}, -2, 0  \right\}~.
\end{equation}
The key feature of this analysis is that the constraints behave differently for $L=(1,2)~\mathrm{mod}~6$ and $L=(3,4,5)~\mathrm{mod}~6$, owing to the contrasting multiplicities of the relevant zeros. This distinction will play an important role in the simplification of the all-level analysis below, and is a recurring feature for $V_k(\sof_{2N})$ with $N$ even.

\subsubsection{Structure of zeros of the vacuum determinant}
It is useful to assemble the smallest conformal level $L_s$ at which $k = k_{2,q}, k_{3,q}$ and $k_{6,q}$ first appear as zeros of the fixed-root vacuum determinant, and the multiplicities of these zeros (which can be computed using \eqref{eq:expinchar}). We document the relevant data in Table \ref{tb:multandLevforso8}.

%%%%%%%%%%%%%%%%%%%%%%%%%%%
\begin{table}[]
\centering
\begin{tabular}{|c|c|c|}
\hline\hline
$k$ & $L_s$ & $m_{p/q}(L_s\delta+\alpha)$ \\
\hline\hline
$k_{2,q}$ & $3q$ & $2$ \\
$k_{3,q}$ & $7q$ & $40$ \\
$k_{6,q}$ & $q$ & $1$ \\
\hline
\end{tabular}
\caption{Smallest conformal level $L_s$, and multiplicities for different values of $k$.}
\label{tb:multandLevforso8}
\end{table}
%%%%%%%%%%%%%%%%%%%%%%%%%%%

Combining this with \eqref{eq:AtconsttypeD}, it follows that the (possible) zeros of the conformal level $L$ vacuum determinant in a fixed-root sector for $V_k(\sof_8)$ can be classified as,
\begin{equation} \label{eq: so8prelimzero}
    k = \begin{cases}
        k_{6,q} = -6 + \frac{6}{q} \ \ : q \leqslant L \ \ \text{and} \ \ (q,6)=1 \\ 
        k_{2,q} = -6 + \frac{2}{q} \ \ : q \leqslant \lfloor \frac{L}{3} \rfloor \ \ \text{and} \ \ (q,2)=1 \\
        k_{3,q} = -6 + \frac{3}{q} \ \ : q \leqslant \lfloor \frac{L}{7} \rfloor \ \ \text{and} \ \ (q,3)=1 \\
        k_{A,B} = -6 + \frac{A}{B}: A \geqslant 8 \ \ \text{and} \ \ B \leqslant L
    \end{cases}
\end{equation}
To implement graded unitarity constraints for $L \geqslant 13$, we first determine some of the smallest zeros of the vacuum determinant in the fixed-root sector. 

Consider $L = (1,2,0)~\textrm{mod}~6$. From \eqref{eq: so8prelimzero} and the data in Table \ref{tb:multandLevforso8}, it follows that for $L = 1~\textrm{mod}~6$, the most negative zero of the fixed root conformal level $L$ vacuum determinant is $k = k_{6,L}$, and the multiplicity of this zero is one. Furthermore, note that for $L \geqslant 13$,
\begin{equation} \label{eq:so8inequality}
    k_{6,L-2} < k_{8,L}~,
\end{equation}
which implies that the second most negative zero of the fixed-root vacuum determinant at conformal level $L$ is $k = k_{6,L-2}$. Similarly, for $L = (2,0)~\textrm{mod}~6$, the most negative zero is $k = k_{6,L-1}$. 

When $L = 3~\textrm{mod}~6$, from Table \ref{tb:multandLevforso8}, it follows that the most negative zero is $k = k_{6,L}$ with multiplicity two, and from \eqref{eq:so8inequality}, the second most negative zero is $k = k_{6,L-2}$. When $L = 4~\textrm{mod}~6$, similar inequality arguments imply that the most negative zero is $k = k_{6,L-1}$ with multiplicity 32, and the second most negative zero is $k = k_{6,L-3}$. Furthermore, for $L = 5~\textrm{mod}~6$, \eqref{eq: so8prelimzero} and the data in Table \ref{tb:multandLevforso8} imply that the most negative zero is $k = k_{6,L}$ with multiplicity one, and the second most negative zero is $k = k_{6,L-2}$ with multiplicity 322.\footnote{It is a technical, but an important point that at $L = 4~\textrm{mod}~6$ and $L = 5~\textrm{mod}~6$, there is no \emph{new} primitive null state at $k = k_{6,L-1}$ and $k = k_{6,L-2}$, respectively, and thus, the evenness of the multiplicity of these zeros follows from the multiplicity of two in the fixed root $L = 3~\textrm{mod}~6$ vacuum determinant. The precise multiplicities of 32 and 322 can be computed by expanding the $1/R$ factor in \eqref{eq:multAKMmain}.} As highlighted by the bespoke analysis, it is relevant to determine the third most negative zero of the vacuum determinant at $L=5~\text{mod}~6$ to implement determinant-based graded unitarity constraints in their full glory. To this end, note that for $L \geqslant 13$, and $L=5~\text{mod}~6$, 
\begin{equation}
    k_{6,L-4} < k_{8,L}~,
\end{equation}
which implies that the third most negative zero of the fixed-root vacuum determinant at $L \geqslant 13$ and $L=5~\text{mod}~6$ is $k_{6,L-4}$. 

To summarise, we obtain the following structure of zeros for the level $L \geqslant 13$ vacuum determinant in a fixed-root sector,
\begin{itemize}
    \item At level $L = 1~\text{mod}~6$, the most negative zero is given by $k=k_{6,L}$, and this zero always occurs with multiplicity one. Moreover, the second most negative zero is given by $k = k_{6,L-2}$. 
    \item At level $L = (0,2)~\text{mod}~6$ and $L$ even, the most negative zero is given by $k=k_{6,L-1}$.
    \item At level $L = 3~\text{mod}~6$, the most negative zero is given by $k = k_{6,L}$, and this zero always occurs with multiplicity two. Moreover, the second most negative zero is given by $k = k_{6,L-2}$. 
    \item At level $L = 4~\text{mod}~6$, the most negative zero is given by $k = k_{6,L-1}$ and this zero always occurs with multiplicity 32. Moreover, the second most negative zero is given by $k = k_{6,L-3}$. 
    \item At level $L = 5~\text{mod}~6$, the most negative zero is given by $k = k_{6,L}$, and this zero always occurs with multiplicity one. Moreover, the second most negative zero is given by $k = k_{6,L-2}$ and this zero always occurs with multiplicity 322. The third most negative zero is given by $k = k_{6,L-4}$.  
\end{itemize}

\subsubsection{Graded unitarity constraints on the level}
The conclusion of the structure of zeros established above is that the level $L \geqslant 13$ vacuum determinant in a fixed-root sector takes the following form,\\
\\
\underline{\textbf{Case 1: } $L = (1,2)~\text{mod}~6$}
\begin{equation} \label{eq:so8vacdet1}
\begin{split}
    &\text{det}~S_{L\delta+\alpha} \propto \textcolor{blue}{(k-k_{6,L})^1}(k-k_{6,L-2})^{\#_1}( \ldots \ldots) \qquad L = 1~\text{mod}~6~, \\
    &\text{det}~S_{L\delta+\alpha} \propto \textcolor{blue}{(k-k_{6,L-1})^{\#_2}}( \ldots \ldots) \qquad \qquad \qquad \ \ L = 2~\text{mod}~6~.
\end{split}
\end{equation}
\\
\\
\underline{\textbf{Case 2: } $L = (3,4,5,0)~\text{mod}~6$}
\begin{equation} \label{eq:so8vacdet2}
\begin{split}
    &\text{det}~S_{L\delta+\alpha} \propto \textcolor{blue}{(k-k_{6,L})^2}(k-k_{6,L-2})^{\#_1}( \ldots \ldots) \qquad \qquad \qquad \qquad \ L = 3~\text{mod}~6~, \\
    &\text{det}~S_{L\delta+\alpha} \propto \textcolor{blue}{(k-k_{6,L-1})^{32}}(k-k_{6,L-3})^{\#_2}( \ldots \ldots) \qquad \qquad \qquad \ \ \ L = 4~\text{mod}~6~, \\
    &\text{det}~S_{L\delta+\alpha} \propto \textcolor{blue}{(k-k_{6,L})^1(k-k_{6,L-2})^{322}}(k-k_{6,L-4})^{\#_3}( \ldots \ldots) \qquad L = 5~\text{mod}~6~, \\
    &\text{det}~S_{L\delta+\alpha} \propto \textcolor{blue}{(k-k_{6,L-1})^{\#_4}}( \ldots \ldots) \qquad \qquad \qquad \qquad \qquad \qquad \ \ \ L = 0~\text{mod}~6~.
\end{split}
\end{equation}

The all-level determinant based graded unitarity constraints now follow in a straightforward manner. In particular, note that when $L = 1~\text{mod}~6$, graded unitarity sign constraints rule out $k \in (k_{6,L},k_{6,L-2})$, whereas when $L = 2~\text{mod}~6$, there are no new constraints. Furthermore, at conformal level $L = (3,4)~\text{mod}~6$, while it is important to note that the most negative zero has changed to $k = k_{6,L}$ or $k = k_{6,L-1}$, there are no new constraints from graded unitarity due to the multiplicity of the most negative zero being even. 

When $L = 5~\text{mod}~6$, graded unitarity sign constraints rule out the possibility of $k \in (k_{6,L}, k_{6,L-2}) \bigcup (k_{6,L-2}, k_{6,L-4})$. On the other hand, when $L = 0~\text{mod}~6$, there are again no new constraints. 

When combined with the low level graded unitarity constraints in \eqref{eq:lowlevelso8fullconst}, we find that with the assumption of the standard $\mfR$-filtration on $V_k(\mathfrak{so}_8)$ affine Kac--Moody algebras, all-level (fixed root) determinant-based graded unitarity constraints rule out all levels other than the following,
\begin{equation} \label{eq:so8alllevel}
    k = \begin{cases}
        \textcolor{knownSCFTtext}{k_{6,q} = -6 + \frac{6}{q}: (6,q) = 1} \\
        \textcolor{knownSCFTtext}{-2} \\
        \textcolor{conjecturedSCFTtext}{-\frac{26}{5}} \\
        k_{2,q} = -6 + \frac{2}{q}: (2,q) = 1 \\
        -\frac{5}{2}
    \end{cases}
\end{equation}
The first set of levels above is the set of boundary admissible levels. The $V_k(\sof_8)$ affine Kac--Moody algebras corresponding to these levels are precisely the ones associated with the ``class-I" $D_q^6(\mathfrak{so}_8)$ Argyres--Douglas SCFTs under the SCFT/VOA correspondence \cite{Xie:2016evu, Song:2017oew}. The level $k = -2$ (while non-admissible) is precisely the one corresponding to the associated VOA $V_{-2}(\sof_8)$ of the Lagrangian $SU(2)$ with $N_f = 4$ flavours SCFT. 

While the level $-\frac{26}{5}$ is non-admissible, $V_{-\frac{26}{5}}(\mathfrak{so}_8)$ has been conjectured to be the associated VOA of the $D^4_5(\mathfrak{so}_8)$ Argyres--Douglas SCFT which has no extra mass parameters/flavour symmetries coming from the irregular puncture. Furthermore, this affine VOA has also been conjectured to be quasi-lisse \cite{adamovic2026new}. Nevertheless, as discussed in Section \ref{sec:intro}, the associated VOA of $D^4_5(\mathfrak{so}_8)$ need not coincide with the affine Kac--Moody algebra $V_{-\frac{26}{5}}(\mathfrak{so}_8)$, but may instead be a non-trivial finite extension thereof. On the other hand, the $V_k(\mathfrak{so}_8)$ affine VOAs corresponding to the remaining set of levels $k_{2,q}$ and  $-\frac{5}{2}$ are not currently known (or conjectured) to be associated with any unitary 4d $\mathcal{N} = 2$ SCFT. In either case, although these levels are not excluded by all-level (fixed root) determinant-based graded unitarity constraints, it remains to be seen whether analogous constraints from a different weight sector or finer state-based constraints exclude any of them. 

\subsection{$V_k(\sof_{2N})$}
We now turn to implementing graded unitarity constraints for $V_k(\sof_{2N})$ with $N$ even. The structure of zeros of the fixed-root vacuum determinant is a natural generalisation of the preceding discussion for $V_k(\mathfrak{so}_8)$.

\subsubsection{Low level graded unitarity constraints}
We begin by performing a low level analysis for conformal levels $ \leqslant 6$. To this end, note that the fixed-root vacuum determinant for $L = 1,2,3,4$ is given by,
\begin{equation} \label{eq: so2NevenL1to3}
\begin{split}
    &\text{det}~S_{\delta + \alpha} \ \propto \textcolor{blue}{(k - 0)}~, \\
    &\text{det}~S_{2\delta + \alpha} \propto \textcolor{blue}{(k - k_{N,1})}(k-k_{2N-2,1})^{\#_1} (\ldots \ldots)~,\\ 
    &\text{det}~S_{3\delta + \alpha} \propto \textcolor{blue}{(k - k_{2N-2,3})}(k-k_{N,1})^{\#_2} (\ldots \ldots)~,\\
    &\text{det}~S_{4\delta + \alpha} \propto \textcolor{blue}{(k - k_{N,2})}(k-k_{2N-2,3})^{\#_3} (\ldots \ldots)~.\\
\end{split}   
\end{equation}
Hence, up to $L = 4$, graded unitarity constraints require the level $k$ to be in the regime,
\begin{equation}
    -(2N-2) < k \leqslant k_{N,2} \qquad \text{or} \qquad k \in \{k_{2N-2,3}, k_{N,1}, 0\}~. 
\end{equation}
At conformal level $L = 5$, one obtains that the fixed-root vacuum determinant is given by,
\begin{equation} \label{eq:level5Dtypeeven}
\begin{split}
    &\text{det}~S_{5\delta + \alpha} \propto \textcolor{blue}{(k - k_{10,5})^2}(k-k_{6,2})^{\#_1} (\ldots \ldots) \qquad \ \ \ \text{when} \ \ N = 6~, \\
    &\text{det}~S_{5\delta + \alpha} \propto \textcolor{blue}{(k - k_{2N-2,5})}(k-k_{N,2})^{\#_2} (\ldots \ldots) \qquad \text{when} \ \ N \geqslant 8~.   
\end{split}
\end{equation}
The contrasting multiplicities of the most negative zero of the fixed root $L = 5$ determinant for $N = 6$ and $N \geqslant 8$ are a straightforward consequence of \eqref{eq:Lmultofk2Nm2}. Hence, there are no new constraints from graded unitarity for $N = 6$, whereas for $N \geqslant 8$, graded unitarity restricts the level $k$ to be in the regime, 
\begin{equation}
    -(2N-2) < k \leqslant k_{2N-2,5} \qquad \text{or} \qquad k \in \{k_{N,2}, k_{2N-2,3}, k_{N,1}, 0\}~. 
\end{equation}
At $L=6$, the fixed-root vacuum determinant is given by,
\begin{equation} \label{eq:Lvl6soeven}
\begin{split}
    &\text{det}~S_{6\delta + \alpha} \ \propto \textcolor{blue}{(k - k_{10,5})^{\mathrm{even}}(k - k_{8,3})^4}(k-k_{6,2})^{\#_1} (\ldots \ldots) \qquad \ \text{when} \ \ N = 6~, \\
    &\text{det}~S_{6\delta + \omega_4} \propto \textcolor{blue}{(k - k_{10,5})^{\mathrm{even}}(k - k_{8,3})^1}(k-k_{6,2})^{\#_2} (\ldots \ldots) \qquad \ \text{when} \ \ N = 6~, \\
    &\text{det}~S_{6\delta + \alpha} \ \propto \textcolor{blue}{(k - k_{N,3})}(k-k_{2N-2,5})^{\#_3} (\ldots \ldots) \qquad \qquad \qquad \ \ \ \text{when} \ \ N \geqslant 8~,   
\end{split}
\end{equation}
where the even multiplicity of $(k-k_{10,5})$ follows from similar arguments as $V_k(\mathfrak{so}_8)$: at $L = (N,N+1)~\mathrm{mod}~(2N-2)$, there is no new primitive null at $k = k_{2N-2,L-1}$ and $k = k_{2N-2,L-2}$, respectively, and the evenness then follows from the multiplicity of two at $L = (N-1)~\mathrm{mod}~(2N-2)$. Note that the multiplicity of four of the $k = k_{8,3}$ zero is a consequence of the null state being a component of the $[0,0,0,1,0,0]$ representation, which has root multiplicity four. Hence, when $N = 6$, there are no new constraints from graded unitarity in the fixed-root vacuum determinant sector. However, considering the vacuum determinant in the weight $\lambda = \omega_4$ sector instead,\footnote{It is worth mentioning that even in the weight $\omega_4$ sector, the multiplicity of $k=k_{6,2}$ is indeed non-zero. This follows from the fact that the primitive null at $k = k_{6,2}$ is a component of the $[0,0,0,0,0,2]$ representation of $\sof_{12}$; and the multiplicity of $\lambda = \omega_4$ in this representation is non-zero. The evenness of the $k = k_{10,5}$ multiplicity persists at $L = 6$ and follows from \eqref{eq:level5Dtypeeven}.} graded unitarity constraints for $N = 6$ eliminate the interval $k \in (k_{8,3}, k_{6,2})$.  

On the other hand, when $N \geqslant 8$, graded unitarity constraints imply,
\begin{equation}
    -(2N-2) < k \leqslant k_{N,3} \qquad \text{or} \qquad k \in \{k_{2N-2,5}, k_{N,2}, k_{2N-2,3}, k_{N,1}, 0\}~. 
\end{equation}
The preceding analysis reveals a systematic alternation between the exclusion of the levels $(k_{2N-2,L}, k_{2N,L-1})$ and $(k_{2N,L}, k_{2N-2,L-1})$. 

\subsubsection{Structure of zeros of the vacuum determinant} \label{sec:soevenstruczero}
Recall that the possible zeros in the conformal level $L$ vacuum determinant are given by \eqref{eq: so2nprelimzero}. The precise structure of zeros in the various regimes of the level $L$ then follows from certain inequalities. For instance, note that for $6 < L \leqslant N+1$ and $N \geqslant 10$,\footnote{When $N = 6$ and $N = 8$, the cases of $L \leqslant 4$ and $L \leqslant 6$ respectively, are already handled by bespoke analysis.} 
\begin{equation} \label{eq:soevenstr1}
    k_{2N, L-1} < k_{4N-8, L}~.
\end{equation}
This implies that for $6 < L \leqslant N-2$ and $N \geqslant 10$, one has the following structure of zeros,
\begin{itemize}
    \item For $L$ odd, the most negative zero of the fixed-root vacuum determinant is $k = k_{2N-2,L}$, and this zero always occurs with multiplicity one. Furthermore, the second most negative zero is $k = k_{2N,L-1}$.
    \item For $L$ even, the most negative zero of the fixed-root vacuum determinant is $k = k_{2N,L}$, and this zero always occurs with multiplicity one. Furthermore, the second most negative zero is $k = k_{2N-2,L-1}$.
\end{itemize}
Consequently, determinant-based constraints from graded unitarity up to $L = N-2$ restrict the level $k$ to be in the regime,
\begin{equation} \label{eq:constSOevenuptoNm2}
    -(2N-2) < k \leqslant k_{2N, N-2} \ \ \text{or} \ \ k \in \{k_{2N-2, N-3}, k_{2N, N-4}, \ldots, k_{2N-2, 3}, k_{2N, 2}, k_{2N-2,1}\}~.
\end{equation}
As a consequence of the subtle nature of the multiplicities of zeros when $L = (N-1)~\text{mod}~(2N-2)$, we will perform a separate analysis for the structure of zeros when $L = (N-1,N,N+1)~\text{mod}~(2N-2)$. To begin with, note that the structure of zeros at $L = N-1, N, N+1$ for $N \geqslant 6$ is given by the following.
\begin{itemize}
    \item At $L = N-1$, the most negative zero of the fixed-root vacuum determinant is given by $k = k_{2N-2, N-1} = k_{2,1}$, and this zero always occurs with multiplicity two. Furthermore, the second most negative zero is $k = k_{2N, N-2}$. 
    \item At $L = N$, 
    \begin{enumerate}
        \item When $N = 6$, the structure of zeros has already been discussed in \eqref{eq:Lvl6soeven}.
        \item On the other hand, for $N \geqslant 8$, the most negative zero of the fixed-root vacuum determinant is given by $k = k_{2,1}$, and this zero always occurs with even multiplicity. Furthermore, the second most negative zero is given by $k = k_{2N, N-2}$. 
    \end{enumerate}
    \item At $L = N+1$, the most negative zero of the fixed-root vacuum determinant is given by $k = k_{2N-2,N+1}$, and this zero always occurs with multiplicity one. 
    \begin{enumerate}
        \item When $N = 6$, the second most negative zero is given by $k = k_{2,1}$, and this zero occurs with even multiplicity. Furthermore, the third most negative zero for $N = 6$ is given by $k = k_{8,3}$. 
        \item On the other hand, for $N \geqslant 8$, the second most negative zero is given by $k = k_{2N-2, N-1} = k_{2,1}$, and this zero occurs with even multiplicity. Moreover, the third most negative zero is given by $k = k_{2N, N-2}$.   
    \end{enumerate}
\end{itemize}
From the above structure of zeros, it follows that up to $L = N+1$, determinant-based graded unitarity constraints rule out all levels $k$ except,
\begin{equation} \label{eq:so4nlevellessNp2}
\begin{split}
    &-10 < k \leqslant k_{10,7} \ \ \text{or} \ \ k \in \{k_{10,i}: i \in \mathcal{I}_6; k_{12,j}: j \in \mathcal{J}_6\} \ \ \text{or} \ \ k = k_{8,3} \qquad \qquad \text{for} \ \ N=6~, \\
    &-(2N-2) < k \leqslant k_{2N-2,N+1} \ \ \text{or} \ \ k \in \{k_{2N-2,i}: i \in \mathcal{I}_N, k_{2N,j}: j \in \mathcal{J}_N \} \qquad \text{for} \ \ N \geqslant 8~, \\
\end{split}
\end{equation}
where $\mathcal{I}_N$ and $\mathcal{J}_N$ are the sets,
\begin{equation}
\begin{split}
    &\mathcal{I}_N = \{1, 3, 5, \ldots, N-1\}~, \\ 
    &\mathcal{J}_N = \{2, 4, 6, \ldots, N-2\}~.
\end{split}
\end{equation}
Now we proceed to describing the structure of zeros for $L \geqslant N+2$, ultimately leading to all-level determinant-based graded unitarity constraints. 

Note that for $N \geqslant 5$ and $L \geqslant N+2$,
\begin{equation}
    k_{2N-2,L-2} < \{k_{2N,L-1}, \ k_{4N-8, L}\}, \qquad k_{2N-2,L-1} < \{k_{2N,L}, \ k_{4N-8, L}\}~,
\end{equation}
and one obtains the following structure of zeros for $L \geqslant N+2$ and $L \neq (N-1,N,N+1)~\text{mod}~(2N-2)$, 
\begin{itemize}
    \item For $L$ odd, the most negative zero is given by $k = k_{2N-2, L}$, and this zero always occurs with multiplicity one. Furthermore, the second most negative zero is given by $k = k_{2N-2, L-2}$. 
    \item For $L$ even, the most negative zero is given by $k = k_{2N-2,L-1}$. 
\end{itemize}
As elaborated previously in this section, having more refined knowledge of the structure of zeros for $L = (N-1,N,N+1)~\text{mod}~(2N-2)$ is desirable. To this end, we note that for $N \geqslant 5$ and $L \geqslant N+2$, the structure of zeros in these cases is given by,
\begin{itemize}
    \item At $L = (N-1)~\text{mod}~(2N-2)$, the most negative zero of the fixed-root vacuum determinant is given by $k = k_{2N-2, L}$, and this zero always occurs with multiplicity two. Furthermore, the second most negative zero is given by $k = k_{2N-2, L-2}$. 
    \item At $L = N~\text{mod}~(2N-2)$, the most negative zero of the fixed-root vacuum determinant is given by $k = k_{2N-2, L-1}$, and this zero always occurs with even multiplicity. Furthermore, one obtains: 
    \begin{enumerate}
        \item For $L = 3N-2$, the second most negative zero is given by $k = k_{2N, 3N-2}$, and this zero occurs with multiplicity one. Furthermore, the third most negative zero is given by $k = k_{2N-2, 3N-5}$.
        \item For $L \neq 3N-2$, the second most negative zero is given by $k = k_{2N-2, L-3}$. 
    \end{enumerate}
    \item At $L = (N+1)~\text{mod}~(2N-2)$, the most negative zero of the fixed-root vacuum determinant is given by $k = k_{2N-2, L}$, and this zero always occurs with multiplicity one. The second most negative zero is given by $k = k_{2N-2,L-2}$, and this zero always occurs with even multiplicity. Furthermore, one obtains: 
    \begin{enumerate}
        \item For $L = 3N-1$, the third most negative zero is given by $k = k_{2N, 3N-2}$.
        \item For $L \neq 3N-1$, the third most negative zero is given by $k = k_{2N-2, L-4}$. 
    \end{enumerate}
\end{itemize}

\subsubsection{Graded unitarity constraints for $V_k(\sof_{2N})$ with $N$ even}
To summarise the structure of zeros of the previous section, one obtains that the fixed-root vacuum determinant for $N \geqslant 6$ and $L \geqslant N+2$ takes the following form, \\
\\
\underline{\textbf{Case 1: } $L = (1,2,\ldots, N-2, N+2, \ldots, 2N-2)~\text{mod}~(2N-2)$}
\begin{equation}
\begin{split}
    &\text{det}~S_{L\delta+\alpha} \propto \textcolor{blue}{(k-k_{2N-2,L})^1}(k-k_{2N-2,L-2})^{\#}( \ldots \ldots) \qquad L \ \ \text{odd}~, \\
    &\text{det}~S_{L\delta+\alpha} \propto \textcolor{blue}{(k-k_{2N-2,L-1})^{\#}}( \ldots \ldots) \qquad \qquad \qquad \qquad \ L \ \ \text{even}~.
\end{split}
\end{equation}

\underline{\textbf{Case 2: } $L = (N-1)~\text{mod}~(2N-2)$}
\begin{equation}
\begin{split}
    &\text{det}~S_{L\delta+\alpha} \propto \textcolor{blue}{(k-k_{2N-2,L})^2}(k-k_{2N-2,L-2})^{\#}( \ldots \ldots)~.
\end{split}
\end{equation}

\underline{\textbf{Case 3: } $L = N~\text{mod}~(2N-2)$}
\begin{equation}
\begin{split}
    &\text{det}~S_{L\delta+\alpha} \propto \textcolor{blue}{(k-k_{2N-2,3N-3})^{\mathrm{even}}(k-k_{2N,3N-2})^1}(k-k_{2N-2, 3N-5})^{\#}( \ldots) \qquad L=3N-2~, \\
    &\text{det}~S_{L\delta+\alpha} \propto \textcolor{blue}{(k-k_{2N-2,L-1})^{\mathrm{even}}}(k - k_{2N-2, L-3})^{\#}( \ldots) \qquad \qquad \qquad \qquad \ \ \ \ \ L \neq 3N-2~.
\end{split}
\end{equation}

\underline{\textbf{Case 4: } $L = (N+1)~\text{mod}~(2N-2)$}
\begin{equation}
\begin{split}
    &\text{det}~S_{L\delta+\alpha} \propto \textcolor{blue}{(k-k_{2N-2,3N-1})^{1}(k-k_{2N-2,3N-3})^{\mathrm{even}}}(k-k_{2N, 3N-2})^{\#}(\ldots) \qquad L=3N-1~, \\
    &\text{det}~S_{L\delta+\alpha} \propto \textcolor{blue}{(k-k_{2N-2,L})^{1} (k - k_{2N-2, L-2})^{\mathrm{even}}}(k - k_{2N-2, L-4})^{\#}(\ldots) \qquad \ \ \ \ L \neq 3N-1~.
\end{split}
\end{equation}

This structure of the fixed-root vacuum determinant makes the implementation of graded unitarity constraints immediate. In particular, note that for $L \geqslant N+2$ and $L = \{1,2,\ldots, N-2, N+2, \ldots, 2N-2\}~\text{mod}~(2N-2)$, determinant-based graded unitarity sign constraints eliminate the region $k \in (k_{2N-2, L}, k_{2N-2, L-2})$ for $L$ odd, whereas, they do not give any constraints for $L$ even. Similarly, for $L = (N-1)~\text{mod}~(2N-2)$ (including $L = 3N-3$), there are no new constraints from graded unitarity. When $L = N~\text{mod}~(2N-2)$, graded unitarity eliminates $k \in (k_{2N, 3N-2}, k_{2N-2, 3N-5})$ for $L = 3N-2$, whereas it imposes no constraints for $L \neq 3N-2$. When $L = (N+1)~\text{mod}~(2N-2)$, graded unitarity constraints further rule out $k \in (k_{2N-2, 3N-1}, k_{2N-2, 3N-3}) \bigcup (k_{2N-2, 3N-3}, k_{2N, 3N-2})$ at $L = 3N-1$, and $k \in (k_{2N-2, L}, k_{2N-2, L-2}) \bigcup (k_{2N-2, L-2}, k_{2N-2, L-4})$ at $L \neq 3N-1$. Combining this with the graded unitarity constraints up to $L = N+1$, \emph{i.e.}, \eqref{eq:so4nlevellessNp2}, one obtains that fixed root determinant-based graded unitarity rules out all levels $k$ except,
\begin{equation} \label{eq:VkSOevengrunitresults}
\begin{split}
    &N = 6: \qquad k = \begin{cases}
        -10 + \frac{10}{q}: (10,q) = 1 \\ 
        -4, -7, -\frac{37}{4} \\
        -10 + \frac{2}{q}:  (2,q) = 1 \\
        -\frac{22}{3}
    \end{cases} \\
    &N \geqslant 8: \qquad k = \begin{cases}
        -(2N-2) + \frac{(2N-2)}{q}: (2N-2, q) = 1 \\
        -(2N-2) + \frac{(2N-2)}{q}: q \in 2\mathbb{N}+1 \ \ \text{and} \ \ (2N-2, q) \neq 1 \\
        -(2N-2) + \frac{N}{q} \qquad : q \in \mathbb{N} \ \ \text{and} \ \ 1 \leqslant q \leqslant \frac{N-2}{2} \\
        -(2N-2) + \frac{2N}{3N-2}
    \end{cases}
\end{split}
\end{equation}
The analysis for $V_k(\mathfrak{so}_{2N})$ with $N$ odd proceeds analogously, although the structure of zeros of the vacuum determinant differs due to the contrasting multiplicities in \eqref{eq:Lmultofk2Nm2}. We defer the details to Appendix \ref{app:VkSOodd}, and combine the resulting restrictions on the level $k$ with those obtained above for even $N$.

We find that with the assumption of the standard $\mfR$-filtration \eqref{eq:AKMRfiltration} on $V_k(\mathfrak{so}_{2N})$ for $N \geqslant 5$, fixed root determinant-based graded unitarity rules out all levels $k$ except,  
\begin{equation} \label{eq:VkSOfullresults}
    k = \begin{cases}
        \textcolor{knownSCFTtext}{-(2N-2) + \frac{(2N-2)}{q}: (2N-2, q) = 1} \\
        \textcolor{conjecturedSCFTtext}{-(2N-2) + \frac{(2N-2)}{q}: q \in 2\mathbb{N}+1 \ \ \text{and} \ \ (2N-2, q) \neq 1} \\
        \textcolor{conjecturedSCFTtext}{-(2N-2) + \frac{N}{q} \qquad : 1 \leqslant q \leqslant \frac{N-2}{2}  \ \ \text{and} \ \  \frac{N}{\text{gcd}(N,q)} \in 2\mathbb{N}} \\
        \textcolor{conjecturedSCFTtext}{-(2N-2) + \frac{2N}{3N-2} \ \ : N \in 4\mathbb{N}} \\
        -(2N-2) + \frac{N}{q} \qquad : 1 \leqslant q \leqslant \frac{N-2}{2} \ \ \text{and} \ \ \frac{N}{\text{gcd}(N,q)} \in 2\mathbb{N}+1  \\
        -(2N-2) + \frac{2N}{3N-2} \ \ : N \in 4\mathbb{N}+2 \\
        -(2N-2) + \frac{2}{q} \qquad \ : q \in 2\mathbb{N}+1 \ \ \text{and} \ \ N \in 2\mathbb{N}+1 \\
        -\frac{22}{3} \qquad \qquad \qquad \ \ \ :N = 6
    \end{cases}
\end{equation}
The first set of levels above is the set of boundary admissible levels, and the $V_k(\sof_{2N})$ affine Kac--Moody algebras corresponding to these levels are precisely the ones associated with the ``class-I" $D_q^{2N-2}(\mathfrak{so}_{2N})$ Argyres--Douglas SCFTs under the SCFT/VOA correspondence \cite{Xie:2016evu, Song:2017oew}. 

While the second, third, and fourth sets of levels represented in green are all non-admissible, the corresponding affine VOAs are conjectured to be associated with the $D_q^{2N-2}(\mathfrak{so}_{2N})$, $D_q^{N}(\mathfrak{so}_{2N})$, and $D_{\frac{3N-2}{2}}^{N}(\mathfrak{so}_{2N})$ Argyres--Douglas SCFTs, respectively.\footnote{It is worth mentioning that while the $V_{-N+2}(\mathfrak{so}_{2N})$ affine Kac--Moody algebras are non-admissible, they are known to be quasi-lisse \cite{arakawa2017sheets}.} As alluded to in Section \ref{sec:intro}, the non-admissibility of these levels means that the associated VOAs of these theories could be non-trivial finite extensions of the corresponding affine Kac--Moody algebras. On the other hand, the remaining set of non-admissible $V_k(\mathfrak{so}_{2N})$ are not currently known to be associated with any unitary 4d $\mathcal{N} = 2$ SCFTs under the SCFT/VOA correspondence. In either case, we leave open the possibility that an analogous analysis in a different weight sector or a finer state-based graded unitarity analysis might rule out some or all of these non-admissible levels.